\section{Auxiliary data and a prime-sum identity}\label{sec:data}

\subsection{Number fields, weights, and levels}

Let \(q_0\) be a positive multiple of \(18480\), and let \(d\) be a prime
congruent to \(q_0/2-1\) modulo \(q_0\). Set \(K=\mathbb Q(\sqrt{-d})\).
Then \(d\equiv7\pmod8\). The character \(\chi_d\) from the preceding section
is the quadratic character of \(K\). We also write it as \(\chi_K\).

The admissible data \(D\) include positive integers \(x,x_1,b\) and integers
\(m,r\) that determine the ranks. Set \(a=xb\), \(k=r+1-m\), and \(t=5r\).
The integers \(m,r\) are positive powers of \(2,3\), respectively, and satisfy
\begin{gather*}
 (\log b)^{100}\le x,\qquad
 x^{10^6}\le x_1,\qquad x_1^{10^6}\le b,\qquad
 (\log x_1)^6\le\log b,\\
 \sqrt{x}\le m\le 2\sqrt{x},\qquad x^5\le r\le 3x^5.
\end{gather*}
Admissibility also requires \(\log\log b\ge
10^{12}\operatorname{constantBound}(q_0,d)\), where the constant and the full
set of conditions are fixed in the formalization. \leanref{Parameters.constantBound}
Write \(2^s=m^{22}\), so that \(s\le32\log x\), and set
\[
  l_0=3^{\lfloor\log_3 s\rfloor+4},\qquad
  q_1=3l_0\,2^s q_0,\qquad q_2=(q_0d)^{s+6},\qquad
  E=K(\zeta_{q_1}).
\]
The integer \(b\) is a multiple of \(q_1x_1\).
Identify the set \(\Phi\) of infinite places of \(E\) with
\((\mathbb Z/q_1\mathbb Z)^\times\). Let \(f=|\Phi|\), and let \(\Phi_a\)
be the set of places congruent to \(1\) modulo \(q_0\).
Choose a subset \(\Phi_b\) of the places congruent to \(-1\) modulo \(q_0\),
and impose the following cardinality and weight conditions.
\[
  f=|\Phi_a|\varphi(q_0),\qquad
  |\Phi_a|\le|\Phi_b|\log b,\qquad
  \kappa(\tau)=
  \begin{cases}
    b=a/x,&\tau\in\Phi_b,\\
    a/x_1,&\tau\notin\Phi_b.
  \end{cases}
\]

The set \(\Phi_b\) satisfies a sparsity condition. For any prime \(\ell\le kt\)
and \(0<i<j<l_0\), there are no \(\sigma,\tau,\rho\in\Phi_b\) satisfying
\(\tau/\sigma\equiv\pm\ell^i\) and \(\rho/\sigma\equiv\pm\ell^j\) simultaneously.
The modulus is the part of \(q_1\) prime to \(\ell\), and the two signs are independent.
The formalization specifies the identification of the infinite places
and the sparsity condition.
\leansrc{New/Reduction/Data/FieldsAndPlaces}{Infinite places}. \leanref{Parameters.IsSparse}

The level primes \(p_0,p_1,N\) are all congruent to \(1\) modulo \(q_2^2\).
The distinct primes \(p_0,p_1\) lie in \([b^{10},b^{11}]\), and \(M=p_0p_1\).
Choose algebraic integers \(\pi_M\) in \(K\) and \(\pi_N\) in \(E\) such that
\(\pi_M\overline{\pi_M}=M\) and \(\pi_N\overline{\pi_N}=N\).
Let \(\bar\kappa=\exp(f^{-1}\sum_{\tau\in\Phi}\log\kappa(\tau))\) be the
geometric mean of the weights and
\(\mathrm{rd}_E=|\operatorname{disc} E|^{1/[E:\mathbb Q]}\) the root discriminant.
The range for \(N\) is
\begin{equation}\label{eq:level-window}
 \log N\in[\log N_-,\log N_-+\sqrt{\log b}],\qquad
 \log N_-=
 m\log\bar\kappa+(m-1)\log\frac{M\mathrm{rd}_E}{2\pi}
 -\frac{15}{2}\log b.
\end{equation}
With this choice of level, the term \(a\log N\) offsets the growth over \(m\)
steps in the archimedean estimate below.

\subsection{Existence of admissible data}

\begin{proposition}\label{prop:admissible-data}
Fix \(q_0,d\) as above and let \(0<\varepsilon<1/100\).
For every sufficiently large integer \(x\), admissible data exist.
One may take \(x_1=\lfloor e^{x^{\varepsilon/6}}\rfloor\), \(m\) to be the
least power of \(2\) at least \(\sqrt{x}\), and \(r\) to be the least power
of \(3\) at least \(x^5\).
The least multiple \(b\) of \(q_1x_1\) exceeding \(e^{x^\varepsilon}\) satisfies
\[
 x^{\varepsilon/6}-1\le\log x_1,\qquad
 x^\varepsilon\le\log b\le x^\varepsilon+1.
\]
\end{proposition}

\begin{proof}
Put \(W=x^{\varepsilon/6}\) and choose \(m,r\) as in the statement.
For \(x_1=\lfloor e^W\rfloor\) one has \(W-1\le\log x_1\le W\).
For all sufficiently large \(x\), \(q_1x_1\le e^{W^6}\), so that the
chosen multiple satisfies \(e^{W^6}<b\le2e^{W^6}\).
It follows that \(\log b\le W^6+1\).
Since \(s=O(\log x)\) and \(W\) is a positive power of \(x\), the required
logarithmic scale conditions hold.

For the sparse set, use Behrend's construction. In an \(n\)-dimensional
box with coordinates in \(\{0,\ldots,H-1\}\), at least \(H^n/(nH^2)\)
points have the same squared norm. Reading each point's coordinates as
digits in base \(2l_0H\) gives a set of exponents.
If \(a_1,a_2\) are positive integers with \(a_1+a_2\le2l_0\), then
\[
 (a_1+a_2)e_2\equiv a_1 e_1+a_2 e_3 \pmod {2^s}
\]
implies \(e_1=e_2=e_3\). Indeed, under \((2l_0H)^n\le2^s\), the integer
representatives on both sides are less than \(2^s\), and the equality
holds digit by digit without carries. The middle point is a strict
weighted average of the endpoints unless they coincide. Since all three
points lie on the same sphere, they coincide. \leanref{exists_three_term_free}

Let \(g=1+3l_0q_0\) and take places of the form \(-g^e\).
If \(2^v\) is the highest power of \(2\) dividing \(q_0\), then \(g\) has
order \(2^s\) modulo \(2^{s+v}\). Hence the selected exponents give
distinct places, all congruent to \(-1\) modulo \(q_0\).
For \(\ell=2\) their ratios are \(1\) modulo \(3l_0\), whereas \(4\) has
order \(l_0\) modulo \(3l_0\). This excludes the forbidden ratios.
For an odd prime \(\ell\), raising the two ratios to \(2j\) and \(2i\),
respectively, and comparing them modulo \(2^{s+v}\) gives
\(2j e_2\equiv2(j-i)e_1+2i e_3\pmod{2^s}\).
The property of the exponents gives \(e_1=e_2=e_3\).
Then \(\ell^{2i}\equiv1\pmod{2^{s+1}}\), which by the order bound in the
formalization would imply \(2^s\le(\ell+1)2i\).
The parameters satisfy \(2l_0(kt+1)<2^s\), yielding a contradiction.
Fix \(n\) in terms of \(\varepsilon\), and choose \(H\) to be a power of \(2\)
with \(H^n\le 2^s/W<(2H)^n\). Under the numerical conditions
\[
 W\le 2^s\le W^n,\qquad (2l_0)^n\le W,\qquad
 3l_0 n2^n\le W^3,
\]
one has \(H\le W\) and the encoded box fits within the modulus.
The number of selected places is at least \(3l_0 2^s/W^6\).
Since \(3l_0 2^s=|\Phi_a|\) and \(W^6\le\log b\), this gives the required
cardinality. \leanref{Parameters.exists_isSparse}

A Siegel--Walfisz type estimate supplies the level primes in windows
of width \(W\) for \(\log p\).
Since \(q_0,d\) are fixed and \(s=O(\log x)\), one can fix \(n\) with
\(q_2^2\le W^{n+7}\).
For \(W\le U\le W^{n+7}\) one can choose a prime in the residue class \(1\)
modulo \(q_2^2\) with \(\log p\) in \([U,U+W]\).
The values of \(U\) used here lie in \([\log b,24m\log b]\).
Let \(h_K\) be the class number of \(K\). One can place \(h_K+1\) disjoint
windows in \([10\log b,11\log b]\).
Each prime \(p\) selected from these windows splits in \(K\). If \(I_p,J_p\)
are the two primes above \(p\), their classes satisfy \([I_p][J_p]=1\).
Choose \(I_{p_0},I_{p_1}\) in the same class. The ideal \(I_{p_0}J_{p_1}\)
is principal, and one may take its generator to be \(\overline{\pi_M}\).
For \(N\), start the window at \(\log N_-\).
Since \(W\le\sqrt{\log b}\), the window lies in the range of
\eqref{eq:level-window}.
Since \(N\equiv1\pmod4\), one can write \(N=u^2+v^2\).
As \(E\) contains \(\sqrt{-1}\), one may take \(\pi_N=u+v\sqrt{-1}\).
The existence argument in the formalization supplies the windows and
verifies the numerical conditions used above. \leanref{eventually_exists_data}
\end{proof}

\subsection{Source and target spaces of the lifts}

The construction uses Hermitian spaces \(X_j,Y_j\) over \(E\).
Their dimensions are \(\dim_E X_j=k+j(t+1)\) and \(\dim_E Y_j=r+jt\),
and \(Y_m\subset X_m\) is a definite subspace of codimension \(1\).
Let \(S_0\) be the space of algebraic automorphic forms on \(X_0\) of the
specified level, with weight \((a,0,\ldots,0)\) at places in \(\Phi_a\)
and weight \(0\) at the remaining infinite places.
The level is principal congruence at the selected places above \(M\).
Modulo \(\pi_N\) the last row is congruent to \((0,\ldots,0,1)\).
The subspace \(S\subseteq S_0\) consists of forms whose average under the
central scalar action supported at \(w,\bar w\) vanishes for every
place \(w\) dividing \(\pi_N\).
\leansrc{New/UpperBound/Objects/Columns}{Source forms and their level}.
\leanref{ramifiedSourceForms}

The principal-level forms on \(Y_0\) carry an action of
\(\bar G=\prod_w\mathrm{GL}_r(\mathbb F_w)\), where \(w\) ranges over the
finite places above \(p_0,p_1\) selected by \(\pi_M\).
Let \(V\) be the image of the projection given by the central idempotent
\(e_0\), and set \(h=\dim_{\mathbb C}V\).
The idempotent \(e_0\) sums the projectors for the irreducible characters
with the following labels. At each \(w\) the label is a monic squarefree
polynomial of degree \(r\) with nonzero constant term, and none of its
roots in the residue field is a \(q_2\)-th root of unity.
Moreover, above each of \(p_0,p_1\) at least one label is a primitive
polynomial, with a root that generates the multiplicative group of the
degree-\(r\) extension.
\leansrc{New/UpperBound/Objects/LabelsAndIndexTypes}{The finite group and its central idempotent}.
\leansrc{New/UpperBound/Objects/Rows}{The target space and its archimedean type}.

Each twist \(\xi\) gives \(m\) successive theta lifts along
\(Y_0\to\cdots\to Y_m\). Their composite restricted to \(V\) is \(\Theta_\xi\).
The character in the \(j\)-th lift is the fixed character \(\mu_j^0\) of
the formalization multiplied by \(\xi^j\).
Write \(1\) for the trivial twist. Applying the normalization and projection of
Section~\ref{sec:upper-bound} to the Poincar\'e map from \(X_0\) to \(Y_0\) gives
\(\mathfrak p_g\colon S\to V\).
\leansrc{New/UpperBound/Objects/RowChain}{The chain of lifts}.
\leansrc{New/Reduction/IntegralTargetForms}{The integral structure for \(\Theta_\xi\)}.

\subsection{Prime sums in the local divisibility estimates}

The image of an integer in \(\Phi\) is its residue class if it is prime
to \(q_1\), and the identity element otherwise. With this convention we
write \(-\nu\tau\). Put \(u_+=\max(u,0)\), and for \(0\le y\le a\) set
\begin{align*}
 s_{\mathrm{fin}}(y,\tau)
   &=\sum_{1\le\nu\le a}\lambda_1(\nu)
      \left(\frac{y}{\nu}-\kappa(-\nu\tau)\right)_+,\\
 s_{\infty}(j)&=\log(j!)+j\log(M/\Omega),\qquad
 \Omega=
    \exp\!\left(\gamma-\RePart\frac{L'}{L}(1,\chi_d)\right)
        \frac{2\pi}{\mathrm{rd}_E},\\
 s_{\mathrm{main}}&=a\bigl(\log(a/\bar\kappa)+c_1\bigr),\qquad
 \mathcal R=\sum_{\tau\in\Phi_a}\bigl(s_{\mathrm{fin}}(a,\tau)-s_{\mathrm{main}}\bigr),\\
 A_0&=\sum_{\tau\in\Phi}
       \bigl(s_\infty(\kappa(\tau))-s_{\mathrm{fin}}(\kappa(\tau),\tau)\bigr),
 \qquad H_0=mkA_0,\qquad E_j=fa(\log b)^j .
\end{align*}
Here \(\gamma\) is Euler's constant and \(c_1=G_1(1)\) is the constant of
\eqref{eq:prime-remainder}.
We will use \(H_0\ge0\), which follows from
\(s_\infty(\kappa(\tau))-s_{\mathrm{fin}}(\kappa(\tau),\tau)
\ge\kappa(\tau)\log b\) at every \(\tau\).
This inequality uses the scale conditions and
\(s_{\mathrm{fin}}(y,\tau)\le4y\log x_1\). \leanref{primeSum_le}

\begin{proposition}\label{prop:total-remainder}
For admissible data,
\begin{equation}\label{eq:total-remainder}
 \mathcal R =
 a\left[
 |\Phi_b|\bigl(R_{q_0}(x)-R_{q_0}(x_1)\bigr)
 +|\Phi_a|R_1(x_1)
 \right].
\end{equation}
\end{proposition}

\begin{proof}
The number of \(\tau\in\Phi_a\) for which \(-\nu\tau\in\Phi_b\) is
\(|\Phi_b|\) when \(\nu\equiv1\pmod{q_0}\), and zero otherwise.
Indeed, if \(\nu\) is prime to \(q_1\), multiplication determines the
image residue class, and \(\Phi_b\) lies in the class \(-1\) modulo \(q_0\).
If \(\nu\) is not prime to \(q_1\), the image of \(-\nu\) is the identity.
Since \(q_1\) and \(q_0\) have the same prime divisors, in this case
\(\nu\) cannot be \(1\) modulo \(q_0\).

Thus, to evaluate \(\sum_{\tau\in\Phi_a}s_{\mathrm{fin}}(a,\tau)\), first
assign the weight \(a/x_1\) at all \(|\Phi_a|\) places. For each term with
\(\nu\equiv1\pmod{q_0}\), change that weight to \(a/x\) at \(|\Phi_b|\) places.
Writing the positive part \(\bigl(a/\nu-a/y\bigr)_+\) as a sum restricted
to \(\nu\le y\) gives
\[
 \frac1a\sum_{\tau\in\Phi_a}s_{\mathrm{fin}}(a,\tau)
  =
 |\Phi_a|B_1(x_1)+|\Phi_b|\bigl(B_{q_0}(x)-B_{q_0}(x_1)\bigr),
 \qquad B_q(y)=R_q(y)+\frac{\log y}{\varphi(q)}+c_q .
\]
The constant \(c_{q_0}\) cancels in the difference. The definition of
the weights also gives
\[
 |\Phi_a|\bigl(\log(a/\bar\kappa)+c_1\bigr)
 =|\Phi_a|(\log x_1+c_1)
   +\frac{|\Phi_b|}{\varphi(q_0)}(\log x-\log x_1).
\]
This is the main term of the preceding sum. Subtracting it yields
\eqref{eq:total-remainder}. \leanref{totalRemainder_eq}
\end{proof}
